%%% FIELD proof-exact-history-compression
Fix \(\delta\in[0,1]\).  We prove the two inclusions by the response-type compression technique.

First let \(P\in\mathfrak M_\delta(b,q,\omega)\), and put \(x_s=P(S=s)\).  Since \(\mathcal S\) is finite, \(x_s\geq0\) and \(\sum_sx_s=1\).  By \(\text{ass:no-vaccine-effect-on-exposure}\), \(P(H^0=H^1)=1\), so \(x_s=0\) on every unequal-history type.  For \(t\in\mathcal T\), \(a\in\mathcal A\), and \(v\in\mathcal V\), randomization, arm positivity, consistency, and the history-indexed response equality give
\[
b_{tav}=P(J_t=v\mid A=a)=P(J_t^a=v)=E\{\widetilde g_{tav}(S)\}=\sum_s\widetilde g_{tav}(s)x_s.
\]
The external-margin transport assumption gives
\[
q_{tv}=P(H_t^0=v)=\sum_s\mathbf 1\{h_t^0=v\}x_s
\]
for every \((t,v)\in\mathcal T\times\mathcal V_0\).  Placebo-reference stability is exactly
\[
\sum_{s,v}\omega_v\{r_{10v}(s)-r^{[0]}_{20v}(s)\}x_s=0.
\]
Finally, for fixed \((a,u,v)\), moving the right side of \(\text{ass:bounded-noninfecting-exposure-priming}\) to the left gives
\[
\sum_s\mathbf 1\{h_1^0=u,r_{1au}=0\}
 \bigl[\mathbf 1\{r^{[u]}_{2av}\ne r^{[0]}_{2av}\}-\delta\bigr]x_s\leq0.
\]
These are precisely all constraints in \(\mathcal P_\delta(b,q,\omega)\), including the cases in which the conditioning-event mass is zero.  Thus every full law projects into the stated polytope.

Conversely, take \(x\in\mathcal P_\delta(b,q,\omega)\).  On a finite probability space draw \(S\) with mass function \(x\), draw \(A\) independently of \(S\) with \(P(A=a)=\pi_a\), and read \(H^0,H^1,R_1,R_2^{[\cdot]}\) from \(S\).  Define \(J^a\) by
\[
\mathbf 1\{J_t^a=v\}=\widetilde g_{tav}(S),\qquad v\in\mathcal V,
\]
and assign \(J_t^a=0\) if all displayed indicators vanish; then set \(J=J^A\).  The definition of \(\widetilde g\) makes at most one variant indicator equal to one in each period, and its period-two survival factor prevents a second first infection, so \(J^a\) is well defined.  This construction gives randomization and consistency by construction, arm positivity from \(\pi_a>0\), and common histories because the unequal-history coordinates of \(x\) vanish.  The observation equalities and the external margins follow from the corresponding equalities defining the polytope.  The placebo equality and every priming inequality follow from its last two constraint families.  Hence the constructed law lies in \(\mathfrak M_\delta(b,q,\omega)\) and reproduces \(b,q,\pi\) and the entire history-indexed response type.  The two constructions are inverse at the level of the type-mass projection, proving equality and the asserted compatible-law certificate.

%%% FIELD proof-extended-sharp-identified-set
Fix \(\delta\) and assume \(\mathcal P_\delta(b,q,\omega)\ne\varnothing\).  For a compatible full law let \(x_s=P(S=s)\).  Theorem \(\text{thm:exact-history-compression}\) gives \(x\in\mathcal P_\delta(b,q,\omega)\), and the controlled risks are, by their definitions,
\[
N(x):=\rho_{11}(x,\omega)=\sum_{s,v}\omega_vr_{11v}(s)x_s,
\qquad
D(x):=\rho_{21}(x,\omega)=\sum_{s,v}\omega_vr^{[0]}_{21v}(s)x_s.
\]
Both belong to \([0,1]\).  Therefore the law's target correspondence is \(\operatorname{Rat}(N(x),D(x))\), and taking the union over compatible laws gives a subset of the three-branch set in \(\text{def:extended-ratio-image}\).

For the reverse inclusion, take any \(x\) appearing in one of those branches.  The converse half of Theorem \(\text{thm:exact-history-compression}\) supplies a full law with exactly this type mass.  If \(D(x)>0\), that law certifies the finite value \(N(x)/D(x)\).  If \(D(x)=0<N(x)\), it certifies \(+\infty\).  If \(N(x)=D(x)=0\), the definition of \(\operatorname{Rat}(0,0)\) makes the same compatible law a certificate for every value of \([0,+\infty]\).  Thus every branch is attainable and the displayed set is the sharp, rather than merely outer, identified set.

%%% FIELD proof-sharp-reciprocal-endpoints
Write \(P=\mathcal P_\delta(b,q,\omega)\), \(N(x)=\rho_{11}(x,\omega)\), and \(D(x)=\rho_{21}(x,\omega)\).  The set \(P\) is a nonempty compact convex polytope: it is the intersection of the probability simplex with finitely many closed affine half-spaces.  By \(\text{ass:reference-risk-positivity}\), \(D(x)\geq\epsilon>0\) throughout \(P\).  Thus the extended branches with zero denominator are absent, and Theorem \(\text{thm:extended-sharp-identified-set}\) yields
\[
\Theta_{I,\delta}(b,q,\omega)=\{N(x)/D(x):x\in P\}.
\]
The ratio is continuous on \(P\).  Since a convex set is connected, its continuous real image is connected and hence an interval; compactness makes that interval closed and makes both extrema attainable.  Also \(0\leq N(x)\leq1\) and \(D(x)\geq\epsilon\), whence the interval is contained in \([0,1/\epsilon]\).

We now verify the reciprocal normalization equation by equation.  Given \(x\in P\), set
\[
z=D(x)^{-1},\qquad y=zx.
\]
Then \(y\geq0\), \(z>0\), \(\sum_sy_s=z\), every homogeneous equality or inequality defining \(P\) is satisfied after multiplication by \(z\), and
\[
\sum_{s,v}\omega_vr^{[0]}_{21v}(s)y_s=zD(x)=1.
\]
Hence \((y,z)\in\mathcal C_\delta(b,q,\omega)\), and its objective is
\[
\sum_{s,v}\omega_vr_{11v}(s)y_s=zN(x)=N(x)/D(x).
\]
Conversely, if \((y,z)\in\mathcal C_\delta(b,q,\omega)\), the normalization equals one, so \(y\ne0\); because \(y\geq0\) and \(\sum_sy_s=z\), necessarily \(z>0\).  Put \(x=y/z\).  Dividing every homogeneous constraint by \(z\) gives \(x\in P\), while
\[
D(x)=z^{-1}\sum_{s,v}\omega_vr^{[0]}_{21v}(s)y_s=z^{-1},
\quad
N(x)/D(x)=\sum_{s,v}\omega_vr_{11v}(s)y_s.
\]
Thus the feasible objective image of the two linear programs is exactly the sharp ratio image.  Its minimizing and maximizing solutions exist because they correspond bijectively to the compact feasible set \(P\) and \(z=1/D(x)\leq1/\epsilon\).  Finally, each optimizing \(x=y/z\) lifts to a compatible endpoint-attaining full law by Theorem \(\text{thm:exact-history-compression}\).

%%% FIELD proof-endpoint-sign-calibration
By Theorem \(\text{thm:sharp-reciprocal-endpoints}\), the compatible values form the attained interval \([L_\delta,U_\delta]\).  An attained interval is contained in \([0,1)\) exactly when its largest element satisfies \(U_\delta<1\), and it is contained in \((1,+\infty]\) exactly when its smallest element satisfies \(L_\delta>1\).  It contains unity exactly when \(L_\delta\leq1\leq U_\delta\).  It contains a value strictly below and a value strictly above unity exactly when \(L_\delta<1<U_\delta\), because both endpoints are attained.  Every endpoint and intermediate value comes from one vector \(x\) carrying both periods, rather than from separately optimized period-specific laws, by the sharp-endpoint theorem.

%%% FIELD proof-common-history-sign-dual
Use the canonical system and abbreviate \(E=E_\omega\), \(F=F_\delta\), \(d=d(b,q)\), and \(c=c_\omega\).  The feasible set is the nonempty compact polytope
\[
P=\{x\geq0:Ex=d,\ Fx\leq0\}.
\]
Consider \(v^+=\max_{x\in P}c^\top x\).  If \(\mu\geq0\) and
\[
E^\top\lambda+F^\top\mu\geq c,
\]
then for every \(x\in P\),
\[
c^\top x\leq\lambda^\top Ex+\mu^\top Fx
\leq d^\top\lambda.
\]
This is weak duality.  We derive the reverse inequality without adding a constraint qualification.  Introduce the finitely generated cone
\[
\mathcal K=\{(Ex,Fx+r,c^\top x-\beta):x\geq0,\ r\geq0,\ \beta\geq0\}.
\]
For a scalar \(t\), the point \((d,0,t)\) belongs to \(\mathcal K\) exactly when some \(x\in P\) has \(c^\top x\geq t\): take \(r=-Fx\) and \(\beta=c^\top x-t\), and conversely read these two equalities from membership.  Hence \((d,0,t)\notin\mathcal K\) for \(t>v^+\).  A finitely generated cone is closed: after normalizing a convergent sequence of conic combinations by the sum of its coefficients, compactness of the convex hull of the finite generators gives either bounded coefficients and a limit representation or an unbounded subsequence whose normalized limit is a zero combination, which can be removed before taking the limit.  The Euclidean projection of \((d,0,t)\) onto this closed convex cone yields a separating vector \((\ell,\nu,\gamma)\) such that
\[
\ell^\top d+\gamma t>0,\qquad
E^\top\ell+F^\top\nu+\gamma c\leq0,\qquad
\nu\leq0,\qquad\gamma\geq0.
\]
Feasibility of \(P\) forces \(\gamma>0\): if \(\gamma=0\), substitute any feasible \(x\) and \(r=-Fx\) into the cone inequality to get \(\ell^\top d\leq0\), contradicting the strict separation.  Therefore
\[
\lambda=-\ell/\gamma,\qquad \mu=-\nu/\gamma
\]
satisfy \(\mu\geq0\), \(E^\top\lambda+F^\top\mu\geq c\), and \(d^\top\lambda<t\).  Letting \(t\downarrow v^+\) proves that the dual infimum is at most \(v^+\), while weak duality gives the reverse inequality.  The dual feasible set is a polyhedron; a finite linear objective with a finite infimum on a nonempty polyhedron attains it, as follows by eliminating any recession direction and minimizing on the resulting bounded face.  Thus a dual optimizer exists with \(d^\top\lambda=v^+\).  Compactness of \(P\) supplies primal attainment and finiteness.

It follows that \(v^+<0\) implies the displayed certificate with \(d^\top\lambda=v^+<0\), while the weak-duality calculation proves the converse.  Moreover,
\[
c^\top x=\sum_{s,v}\omega_v\{r_{11v}(s)-r^{[0]}_{21v}(s)\}x_s
=\rho_{11}(x,\omega)-\rho_{21}(x,\omega).
\]
Thus \(v^+<0\) forces a positive denominator and a ratio below one for every compatible \(x\).

For the other side let \(v^-=\min_{x\in P}c^\top x\).  The same separation argument applied to \(-c\), or direct weak duality, gives
\[
v^-=\max_{\eta,\kappa}\{d^\top\eta:
E^\top\eta-F^\top\kappa\leq c,\ \kappa\geq0\}.
\]
Indeed, for every feasible \((\eta,\kappa)\) and \(x\in P\),
\[
c^\top x\geq\eta^\top Ex-\kappa^\top Fx
\geq d^\top\eta.
\]
Consequently \(v^->0\) is equivalent to the stated certificate with \(d^\top\eta>0\), and it forces every extended ratio value above one (if the denominator is zero, the positive contrast makes the numerator positive and the value is \(+\infty\)).

If the below-unity certificate fails, strong duality implies \(v^+\geq0\); an attaining maximizer \(x^+\) therefore has a ratio value on or above unity, with the \(0/0\) convention supplying such a value when necessary.  If the above-unity certificate fails, an attaining minimizer \(x^-\) has \(c^\top x^-\leq0\) and supplies a value on or below unity.  Theorem \(\text{thm:exact-history-compression}\) lifts either optimizer to a compatible full law, proving the counterlaw assertion.

%%% FIELD proof-biological-waning-translation
Put \(p=\rho_{10}(x,\omega)=\rho_{20}(x,\omega)>0\).  The assumed fiber-wide reference-risk positivity gives \(\rho_{21}(x,\omega)\geq\epsilon>0\), so \(\Psi_\omega(x)\) is the singleton ratio.  Direct algebra gives
\[
\operatorname{VE}_1(x)-\operatorname{VE}_2(x)
=\left(1-\frac{\rho_{11}(x,\omega)}p\right)
-\left(1-\frac{\rho_{21}(x,\omega)}p\right)
=\frac{\rho_{21}(x,\omega)-\rho_{11}(x,\omega)}p.
\]
Since \(p>0\), this difference is positive exactly when \(\rho_{11}/\rho_{21}<1\), and negative exactly when \(\rho_{11}/\rho_{21}>1\).  The common placebo denominator is the extra condition that permits the ratio sign to be read as a temporal change in biological vaccine effect; endpoint sign calibration supplies the corresponding identified-set language.

%%% FIELD proof-k2-worked-compatible-lift
The four atoms have equal mass and common history \((i,j)\), so each period's exposure margin is uniform on variants one and two and has zero mass on no exposure.  In arm zero, period-one infection occurs exactly on the two atoms with \(i=1\), giving \(b^{(2)}_{1,0,1}=1/2\).  Period-two infection occurs exactly when \(i=2\) (period-one survival) and \(j=1\), namely on \(s_{21}\), giving \(b^{(2)}_{2,0,1}=1/4\).  In arm one, the analogous events are \(i=2\) in period one and \((i,j)=(1,2)\) in period two, giving \(b^{(2)}_{1,1,2}=1/2\) and \(b^{(2)}_{2,1,2}=1/4\).  Every other trial cell is zero.

For all \(u\in\mathcal V_0\), the stipulated period-two response equals its \(u=0\) counterpart.  Hence every priming-disagreement indicator is zero, so every sensitivity inequality holds for every \(\delta\in[0,1]\).  The common-history support, probability, trial-risk, and exposure equalities have just been checked.  Under the half--half reference mixture, arm zero responds only to variant one and arm one only to variant two, in both periods.  Therefore, for \(a=0,1\) and \(t=1,2\),
\[
\rho_{ta}(x_\star^{(2)},\omega_\star^{(2)})=\tfrac12,
\qquad
\Psi_{\omega_\star^{(2)}}(x_\star^{(2)})=\{(1/2)/(1/2)\}=\{1\}.
\]
In particular the placebo-stability equality holds.  Thus all defining constraints of the reduced polytope hold, and the explicit construction in \(\text{def:k2-worked-lift}\), equivalently the converse lift in Theorem \(\text{thm:exact-history-compression}\), is a member of the asserted full-law class.

%%% FIELD proof-single-variant-reduction
Apply the displayed constructions to their one-variant boundary specialization.  The exposure restrictions \(q_{11}=q_{21}=1\) and common histories imply \(H_1^0=H_2^0=1\) almost surely.  Write \(R_1=R_{1a1}\), \(R_2^0=R^{[0]}_{2a1}\), and \(R_2^1=R^{[1]}_{2a1}\).  The \(\delta=0\) priming constraint is
\[
P(R_1=0,R_2^1\ne R_2^0)=0.
\]
Consequently the observation map and the trial equalities yield
\[
b_{1a1}=E(R_1),
\qquad
b_{2a1}=E\{R_2^1(1-R_1)\}=E\{R_2^0(1-R_1)\}.
\]
Since \(\omega_1=1\), \(\rho_{2a}(x,\omega)=E(R_2^0)\), and hence
\[
\rho_{2a}(x,\omega)
=b_{2a1}+E(R_2^0R_1).
\]
The binary support gives \(0\leq R_2^0R_1\leq R_1\).  Taking expectations and using \(E(R_1)=b_{1a1}\) proves
\[
b_{2a1}\leq\rho_{2a}(x,\omega)\leq b_{1a1}+b_{2a1}.
\]

%%% FIELD proof-sensitivity-range-report
For each feasible \(\delta\in\mathcal D\), Theorem \(\text{thm:extended-sharp-identified-set}\) identifies exactly the target values generated by \(\mathfrak M_\delta(b,q,\omega)\).  Taking the union over \(\delta\) therefore gives exactly the values generated by the union of full-law classes.  Its closure is, by definition, the smallest closed identified set containing all such values, namely \(\Theta_I^{\mathcal D}(b,q,\omega)\).  Every point before closure has a full-law certificate, and every point added by closure is a limit of such certified values; this is the asserted sharp closed, rather than merely outer, envelope.

For a feasible \(\delta\), write
\[
M(\delta)=\max_{x\in\mathcal P_\delta(b,q,\omega)}c_\omega^\top x,
\]
with \(M(\delta)=-\infty\) when the fiber is empty.  Theorem \(\text{thm:common-history-sign-dual}\) says pointwise that all compatible values are below unity exactly when \(M(\delta)<0\).  A uniform strict conclusion means that one margin \(\gamma>0\) works for every reported feasible \(\delta\): \(c_\omega^\top x\leq-\gamma\) for all such \((\delta,x)\).  This is equivalent, equation by equation, to
\[
\sup_{\delta\in\mathcal D}M(\delta)\leq-\gamma<0,
\]
and hence to the criterion in the statement.  Replacing maxima of \(c_\omega^\top x\) by minima and reversing the inequalities gives the uniform above-unity criterion.  The prescribed default \(\mathcal D=[0,1]\) follows directly from the definition of the sensitivity range when no controlled validation narrows it.

%%% FIELD statement-case-margin-invariance
Fix \(\delta\in\mathcal D\) and assume the case-frequency denominators are positive.  Put \(q^0=q\), \(q^1=q^{\mathrm{case}}\), \(P_j=\mathcal P_\delta(b,q^j,\omega)\), \(N(x)=\rho_{11}(x,\omega)\), and \(D(x)=\rho_{21}(x,\omega)\).  There is a finite necessary-and-sufficient observable primal--dual test for endpoint invariance.  For each \(j\), first test feasibility of \(P_j\), of \(P_j\cap\{N=D=0\}\), of \(P_j\cap\{D=0\}\), and of \(P_j\cap\{D>0\}\) by linear programs.  If \(P_j\cap\{N=D=0\}\ne\varnothing\), its endpoint code is \((0,+\infty)\).  If that set is empty but \(D=0\) on all of \(P_j\), its code is \((+\infty,+\infty)\).  Otherwise form the Charnes--Cooper program with \(q^j\): minimize and maximize \(N(y)\) subject to \(y,z\geq0\), \(E_\omega y=d(b,q^j)z\), \(F_\delta y\leq0\), and \(D(y)=1\).  Its minimum is \(L_j\); its maximum is \(U_j\) when \(P_j\cap\{D=0\}=\varnothing\), and \(U_j=+\infty\) otherwise.  Every finite endpoint is certified by an attaining primal pair and a matching dual pair.  The replacement \(q\mapsto q^{\mathrm{case}}\) leaves both sharp endpoints unchanged if and only if \(P_0,P_1\ne\varnothing\) and their two endpoint codes agree componentwise.  These conditions depend only on \((b,q,q^{\mathrm{case}},\omega,\delta)\).

If invariance fails, the same finite programs give the exhaustive requested certificate.  If \(P_1=\varnothing\), there are \(\lambda\) and \(\mu\geq0\) with \(E_\omega^\top\lambda+F_\delta^\top\mu\geq0\) and \(d(b,q^{\mathrm{case}})^\top\lambda<0\).  If \(P_1\ne\varnothing\) and its extended ratio correspondence is nonconstant, two feasible type masses, with target selections when a \(0/0\) type is used, lift to compatible full laws with different ratios; a primal optimizer and matching dual value for an endpoint whose code differs from that for \(P_0\) certify strict mismatch.  If \(P_1\ne\varnothing\) and its ratio correspondence is constant, a feasible type mass lifts to a compatible law attaining that singleton, and matching lower and upper dual values show that this singleton differs from the original sharp set.  Thus case-strain frequencies are valid substitutes exactly on this computable invariance frontier.

%%% FIELD proof-case-margin-invariance
All programs are finite because \(\overline{\mathcal S}\) is finite.  Write \(E=E_\omega\), \(F=F_\delta\), \(d_j=d(b,q^j)\), and let \(n,d_2\) be the coefficient vectors of \(N,D\).  Feasibility of \(P_j=\{x\geq0:Ex=d_j,Fx\leq0\}\) is an ordinary linear program.  Its finite alternative is
\[
P_j=\varnothing
\quad\Longleftrightarrow\quad
\exists(\lambda,\mu):\ \mu\geq0,\quad E^\top\lambda+F^\top\mu\geq0,
\quad d_j^\top\lambda<0.
\]
The forward implication is obtained by separating \(d_j\) from the finitely generated feasible cone; the reverse follows by multiplying the componentwise inequality by any putative \(x\), which would give \(0\leq d_j^\top\lambda+\mu^\top Fx\leq d_j^\top\lambda<0\).  Adding \(n^\top x=d_2^\top x=0\), or \(d_2^\top x=0\), gives the two zero-risk tests.  Existence of \(D>0\) is certified by \(\max_{x\in P_j}d_2^\top x>0\).  Hence all branch tests are closed finite linear programs.

Suppose first that a \(0/0\) type mass is feasible.  By the definition of \(\operatorname{Rat}(0,0)\), every extended value is compatible, so the code is \((0,+\infty)\).  If no such mass exists and \(D=0\) everywhere, compactness gives \(\min_{P_j}N>0\); every value is therefore \(+\infty\).  In the remaining case there is a point with \(D>0\).  The transformation \(y=x/D(x),z=1/D(x)\) and its inverse \(x=y/z\) prove exactly as in Theorem \(\text{thm:sharp-reciprocal-endpoints}\) that the finite-ratio image equals the objective image of the stated normalization.  If some feasible \(x\) has \(D=0<N\), ratios along a segment from a positive-denominator point to that point diverge, and \(+\infty\) is itself compatible, so \(U_j=+\infty\).  If no such point exists, compactness makes \(D\) uniformly positive and both normalized extrema are finite and attained.  If a minimizing sequence in the normalized program escaped with \(z\to\infty\), its corresponding \(x\)'s would converge to a zero-denominator point; absence of \(0/0\) would make \(N/D\to+\infty\), contradicting minimization.  Thus the finite lower endpoint is attained as well.

For completeness, the minimizing normalization has the observable dual
\[
\max_{\lambda,\kappa,\tau}\ \tau
\quad\text{subject to}\quad
E^\top\lambda-F^\top\kappa+\tau d_2\leq n,
\quad d_j^\top\lambda\geq0,\quad\kappa\geq0,
\]
and the maximizing normalization has the dual
\[
\min_{\lambda,\mu,\sigma}\ \sigma
\quad\text{subject to}\quad
E^\top\lambda+F^\top\mu+\sigma d_2\geq n,
\quad d_j^\top\lambda\leq0,\quad\mu\geq0.
\]
Multiplication by feasible \((y,z)\) proves weak duality in each display.  Separation of the finite polyhedral constraint cone, as in the proof of Theorem \(\text{thm:common-history-sign-dual}\), proves equality whenever the corresponding endpoint is finite.  Thus equality of the two branch codes is both necessary and sufficient and is witnessed entirely by finitely many primal and dual solutions.

If invariance fails, at least one endpoint code differs.  The infeasible substituted-fiber branch has the displayed Farkas ray.  In the nonempty nonconstant branch, the definition of nonconstancy supplies two type masses (or two target selections on a \(0/0\) mass) with different extended values; exact history compression lifts them.  The mismatched code selects a substituted endpoint witness, while its matching dual gives an equality bound that excludes equality with the original endpoint.  In the constant branch both substituted endpoints equal its single value; one lifted feasible mass attains it and matching lower and upper certificates pin it down.  Because invariance failed, the original endpoint pair is different.  These three cases are disjoint and exhaustive.

%%% FIELD statement-perenyi-collapse-boundary
Let \(K=2\), fix \(t\in\mathcal T\) and \(\delta\in\mathcal D\), and suppose the observed at-risk denominators and the four observed risks needed for \(\operatorname{CSE}^{\mathrm{obs}}_t(b)=k_{\mathrm{obs}}>0\) are positive.  Let \(a,c,d,e\) be the coefficient vectors on \(\overline{\mathcal S}\) of, respectively, \(\rho_{t1}(\cdot,e^{(1)})\), \(\rho_{t0}(\cdot,e^{(2)})\), \(\rho_{t0}(\cdot,e^{(1)})\), and \(\rho_{t1}(\cdot,e^{(2)})\), and let \(F^+=\{x\in\mathcal P_\delta(b,q,\omega):a^\top x,c^\top x,d^\top x,e^\top x>0\}\).  The face is nonempty exactly when the linear program maximizing \(\tau\) subject to \(x\in\mathcal P_\delta\) and each of the four risks at least \(\tau\) has value \(\tau^*>0\).  If it is nonempty, choose \(x^\circ\in\operatorname{relint}(\mathcal P_\delta)\cap F^+\), and let the columns of \(B\) span \(\operatorname{lin}(\mathcal P_\delta-\mathcal P_\delta)\).  Define
\[
H_k=\tfrac12\{ac^\top+ca^\top-k(de^\top+ed^\top)\}.
\]
Then \(\operatorname{CSE}^{\mathrm P}_t(x)\) is constant on \(F^+\) and equals \(\operatorname{CSE}^{\mathrm{obs}}_t(b)\) if and only if
\[
(x^\circ)^\top H_{k_{\mathrm{obs}}}x^\circ=0,
\qquad B^\top H_{k_{\mathrm{obs}}}x^\circ=0,
\qquad B^\top H_{k_{\mathrm{obs}}}B=0.
\]
All entries in this test are obtained by finite linear algebra and linear programming from \((b,q,\omega,\delta)\).

Failure has an exhaustive finite certificate.  Empty \(F^+\) is certified by a dual optimum \(\tau^*=0\) (or by infeasibility of the base polytope).  On a nonempty face, put \(k^\circ=\operatorname{CSE}^{\mathrm P}_t(x^\circ)\).  If the same three identities fail with \(k^\circ\), an explicit sufficiently small affine perturbation \(x^\circ+Bh\) gives a second feasible point with a different value, and both points lift to compatible full laws.  If those identities hold with \(k^\circ\) but \(k^\circ\ne k_{\mathrm{obs}}\), the lift of \(x^\circ\), together with its four controlled-risk calculations, is an observable-mismatch certificate.  The same empty-face certificate applies when an observed positivity condition required to define \(k_{\mathrm{obs}}\) fails.

%%% FIELD proof-perenyi-collapse-boundary
On \(F^+\), direct cancellation of the four positive component denominators gives
\[
\operatorname{CSE}^{\mathrm P}_t(x)
=\frac{(a^\top x)(c^\top x)}{(d^\top x)(e^\top x)}.
\]
The phase-I program has \(\tau^*>0\) exactly when one feasible \(x\) makes all four risks positive: one direction is immediate, while the other uses the attained positive minimum of four positive numbers.  Its constraint set is a finite polytope, so the primal--dual separation argument gives a finite matching dual certificate; if the base polytope is empty, the Farkas ray from the preceding theorem applies.

Assume now that \(F^+\ne\varnothing\).  A convex combination of any point in \(F^+\) with a relative-interior point of the polytope lies in \(\operatorname{relint}(\mathcal P_\delta)\cap F^+\), so \(x^\circ\) exists and is computable by a finite phase-I program.  For any \(k>0\), define
\[
Q_k(x)=(a^\top x)(c^\top x)-k(d^\top x)(e^\top x)=x^\top H_kx.
\]
Every point in the affine hull is \(x^\circ+Bh\).  Expanding the quadratic gives
\[
Q_k(x^\circ+Bh)
=(x^\circ)^\top H_kx^\circ
+2h^\top B^\top H_kx^\circ+h^\top B^\top H_kBh.
\]
The polynomial vanishes for every \(h\) exactly when its constant, linear, and quadratic coefficients are all zero, which is precisely the three displayed matrix identities.  Because \(x^\circ\) is relative interior and all four risks there are positive, a neighborhood of \(h=0\) maps into \(F^+\).  Hence a quadratic polynomial vanishing on \(F^+\) vanishes on an open neighborhood and therefore has all three coefficients zero; conversely those identities make it vanish throughout the affine hull.  Since the denominator in the displayed subtype-effect formula is positive on \(F^+\), \(Q_{k_{\mathrm{obs}}}=0\) is equivalent to equality with the observed functional.  This proves necessity and sufficiency.

For the failure partition, \(k^\circ\) makes \(Q_{k^\circ}(x^\circ)=0\).  If a remaining coefficient is nonzero, choose \(h\) for which the resulting one-variable polynomial in a sufficiently small scalar multiple of \(h\) is nonzero.  Relative interior and strict positivity keep that perturbed point in \(F^+\), and \(Q_{k^\circ}\ne0\) makes its ratio differ from that at \(x^\circ\).  Exact history compression lifts both points.  If every coefficient vanishes, the ratio is constantly \(k^\circ\); failure of the test with \(k_{\mathrm{obs}}\) then means \(k^\circ\ne k_{\mathrm{obs}}\), and the four displayed linear forms at \(x^\circ\) give the promised mismatch calculation.  These cases, together with emptiness, are exhaustive.

%%% FIELD statement-mixture-sign-radius
Fix a reported \(\delta\) with nonempty fixed fiber \(P=\mathcal P_\delta(b,q,\omega)\) and one strict baseline sign.  Put \(a_v(x)=\sum_s\{r_{11v}(s)-r^{[0]}_{21v}(s)\}x_s\), and set \(\sigma=1\) when \(U_\delta<1\) and \(\sigma=-1\) when \(L_\delta>1\).  Let \(x^1,\ldots,x^M\) be the vertices of \(P\).  For each \(i\), solve the finite linear program
\[
r_i=\min_{w,t}\sum_vt_v
\quad\text{subject to}\quad w\in\Delta(\mathcal V),\quad
t_v\geq w_v-\omega_v,\quad t_v\geq\omega_v-w_v,\quad
\sigma\sum_vw_va_v(x^i)\geq0,
\]
with value \(+\infty\) if infeasible.  Then \(r_\star(\delta)=\min_i r_i\).  If finite, a minimizing vertex and mixture attain the radius and satisfy the boundary equality \(\sum_vw_va_v(x^i)=0\); the vertex lifts to a compatible full law in the fixed \((b,q,\delta,\omega)\)-fiber.

For any \(r<r_\star(\delta)\), enumerate the vertices \(w^1,\ldots,w^J\) of \(W_r=\{w\in\Delta(\mathcal V):\lVert w-\omega\rVert_1\leq r\}\).  For every \(j\) there are \(\lambda_j\) and \(\mu_j\geq0\) satisfying
\[
E_\omega^\top\lambda_j+F_\delta^\top\mu_j\geq \sigma c_{w^j},
\qquad d(b,q)^\top\lambda_j<0,
\]
where \(c_w(s)=\sum_vw_v\{r_{11v}(s)-r^{[0]}_{21v}(s)\}\).  This finite collection is a robust dual certificate that every \(x\in P\) and every \(w\in W_r\) retain the baseline strict side.  Finally, \(r_\star(\delta)=+\infty\) exactly when all the vertex programs are infeasible; equivalently, strict-sign dual certificates exist for every point-mass mixture \(e^{(v)}\).  These certificates prove infeasibility of reversal for every mixture.  Throughout, only the target mixture changes; the compatible-law fiber and its placebo row remain indexed by the original \(\omega\).

%%% FIELD proof-mixture-sign-radius
Endpoint sign calibration gives \(\sigma c_\omega^\top x<0\) for every \(x\in P\), where \(c_w^\top x=\sum_vw_va_v(x)\).  For a fixed candidate mixture \(w\), the maximum of the linear functional \(\sigma c_w^\top x\) over the compact polytope \(P\) occurs at a vertex.  Consequently
\[
\exists x\in P:\sigma c_w^\top x\geq0
\quad\Longleftrightarrow\quad
\exists i:\sigma\sum_vw_va_v(x^i)\geq0.
\]
Minimizing \(\lVert w-\omega\rVert_1\) over this finite union gives \(\min_i r_i\), and the auxiliary variables \(t_v\) linearize the norm exactly.  This proves the formula for \(r_\star\).  Compactness of the simplex gives attainment whenever one program is feasible.  At a minimizing pair the contrast must equal zero: if it were strictly positive, the segment from \(w\) toward \(\omega\) would, by continuity and the strict negative baseline contrast at the same \(x^i\), meet zero at a strictly smaller distance.

Now fix \(r<r_\star\).  For each vertex \(w^j\) of \(W_r\), the preceding equivalence gives
\[
\max_{x\in P}\sigma c_{w^j}^\top x<0.
\]
The finite polyhedral dual derivation in Theorem \(\text{thm:common-history-sign-dual}\) therefore produces the displayed \((\lambda_j,\mu_j)\).  Any \(w\in W_r\) is a convex combination \(w=\sum_j\gamma_jw^j\).  For every \(x\in P\), weak duality and linearity in \(w\) give
\[
\sigma c_w^\top x
=\sum_j\gamma_j\sigma c_{w^j}^\top x
\leq\sum_j\gamma_jd(b,q)^\top\lambda_j<0.
\]
Thus the finite family certifies the strict side simultaneously over the whole product set.

If every vertex reversal program is infeasible, no pair \((x,w)\in P\times\Delta(\mathcal V)\) reverses the sign, so the defining minimum is over the empty set and equals \(+\infty\).  Conversely, a feasible program gives a finite radius.  Since the simplex is the convex hull of its point masses, dual strict-sign certificates for every \(e^{(v)}\) imply the sign for every mixture by the same convex-combination calculation, and necessity follows by applying fixed-mixture strong duality.  The compatible \(x\) is never moved to a fiber indexed by \(w\), which proves the fixed-fiber assertion.

%%% FIELD statement-exact-whole-set-projection
At a prespecified \(\delta\), under the stated conditional multinomial trial model and period-marginal multinomial external sampling model,
\[
\Pr\{\Theta_{I,\delta}(b,q,\omega)\subseteq
\mathcal C_{n,m,\delta}^{\mathrm{exact}}\mid\text{realized arm sizes}\}\geq1-\alpha
\]
for every compatible full law, including all zero-cell and extended zero-denominator cases.  A finite optimizer-free representation is obtained with \(s=\theta/(1+\theta)\in[0,1]\), where \(s=1\) denotes \(+\infty\): the confidence set is the inverse image of
\[
\left\{s\in[0,1]:\exists(b',q')\in\mathcal G_{n,m}(\alpha),\ exists x\in\mathcal P_\delta(b',q',\omega),\
(1-s)\rho_{11}(x,\omega)=s\rho_{21}(x,\omega)\right\}.
\]
Each exact multinomial acceptance region has a finite support-stratified polynomial description, so the display is the projection of a finite union of compact semialgebraic sets.  Candidate pairs with empty fibers contribute no point, no endpoint optimizer is selected, \(s=1\) represents every compatible zero-denominator branch, and a \(0/0\) law satisfies the equation for every \(s\).

%%% FIELD proof-exact-whole-set-projection
We first prove exact validity of one inversion.  For a fixed multinomial parameter \(p\), the count sample space is finite.  Write \(w_p(u)=\operatorname{Mult}_N(u;p)\) and
\[
\operatorname{pv}_N(p;u)=\sum_{z:T_p(z)\geq T_p(u)}w_p(z).
\]
Let \(R_\gamma=\{u:\operatorname{pv}_N(p;u)<\gamma\}\).  If this set is nonempty, choose \(u_*\in R_\gamma\) with the smallest value of \(T_p\).  Then \(R_\gamma\subseteq\{z:T_p(z)\geq T_p(u_*)\}\), whence
\[
P_p(R_\gamma)\leq\operatorname{pv}_N(p;u_*)<\gamma.
\]
If it is empty the inequality is immediate.  Thus \(P_p\{p\in\Gamma_N(U;\gamma)\}\geq1-\gamma\).  This argument uses the finite ordering only and remains valid for \(N=0\), zero observed cells, and boundary candidate probabilities, with the conventions in the definition.

Conditional on the realized treatment-arm sizes, apply this result to the two arm multinomials at level \(\alpha/4\).  Apply it also to each of the two period-specific external multinomial margins.  Independence between the two external period counts is unnecessary.  The union bound gives
\[
P\{(b,q)\in\mathcal G_{n,m}(\alpha)\mid\text{arm sizes}\}\geq1-4(\alpha/4)=1-\alpha.
\]
For any full law in \(\mathfrak M_\delta(b,q,\omega)\), exact history compression makes \(\mathcal P_\delta(b,q,\omega)\) nonempty.  On the event \((b,q)\in\mathcal G_{n,m}(\alpha)\), the union defining \(\mathcal C_{n,m,\delta}^{\mathrm{exact}}\) includes the entire true set \(\Theta_{I,\delta}(b,q,\omega)\); closure can only enlarge it.  This proves simultaneous whole-set coverage, including \(+\infty\) and \(0/0\) branches because no division was used.

It remains to justify the finite representation.  For each fixed total \(N\) there are finitely many count vectors.  Stratify the simplex by its finitely many support faces and by all possible weak orderings of the finitely many values \(T_p(z)\).  On a positive support face, exponentiating a comparison \(T_p(z)\geq T_p(u)\) and clearing the fixed integer powers writes it as a polynomial monomial inequality; zero-support cases are separate faces.  On each ordering cell the tail probability is a finite sum of multinomial monomials, so \(\operatorname{pv}_N(p;u)\geq\gamma\) is a polynomial inequality.  Hence every \(\Gamma_N\), and their finite intersection \(\mathcal G_{n,m}(\alpha)\), has a finite support-stratified semialgebraic description.  The weak-order convention includes ties and makes the acceptance set closed; it is therefore compact inside the product of simplices.

Finally the map \([0,+\infty]\to[0,1]\), \(\theta\mapsto s=\theta/(1+\theta)\), with \(+\infty\mapsto1\), is a homeomorphism.  For nonnegative \(N,D\),
\[
(1-s)N=sD
\]
is equivalent to \(\theta=N/D\) when \(D>0\); at \(s=1\) it is equivalent to \(D=0\); and when \(N=D=0\) it holds for every \(s\).  Thus it reproduces exactly the three branches of \(\operatorname{Rat}\).  The candidate graph, the reduced-polytope constraints, and this equation form a compact finite semialgebraic set, whose projection is compact and hence already closed.  Empty fibers have no existential witness, and the construction quantifies over all \(x\) rather than selecting any optimizer.
